\documentclass[10pt,letterpaper]{article}
\usepackage[top=1in,left=1in,right=1in,bottom=1in]{geometry}
\usepackage{microtype}
\usepackage{booktabs}
\usepackage{xcolor}
\usepackage[round,authoryear]{natbib}
\usepackage[colorlinks=true,allcolors=blue]{hyperref}

\setlength{\parindent}{0pt}
\setlength{\parskip}{1em}

\title{\vspace*{-0.8in}\textbf{Response to Reviewers --- Fourth Revision}}
\author{}
\date{}

\begin{document}
\maketitle

\noindent\textbf{Submission ID:} fbf130b1-76af-4698-aa61-7b55e8721359\\
\textbf{Manuscript:} A Four-Dimension Pre-Modeling Audit Protocol for Educational Prediction Benchmarks\\
\textbf{Journal:} Scientific Reports\\
\textbf{Revision:} Fourth round

\vspace{0.2in}
\noindent Dear Editor and Reviewers,

We thank the editor and reviewers for their continued engagement with our
manuscript. We are grateful to Reviewer~3 for the positive assessment
and for confirming that the paper is now ready for acceptance; no
further issues were raised by Reviewer~3 in this round. Reviewer~2
raised two further issues, both concerning the interpretation and scope
of the feature-attribution analysis (Section~3.6). Both points are
well-taken: the previous wording overstated the strength of the
conclusions and generalized the feature-attribution findings beyond the
three datasets actually analysed. We have revised the feature-attribution
section accordingly, softening the interpretive claims to exploratory
associations that are consistent with---but do not validate---the audit
profiles, and restricting descriptive statements to the three analysed
datasets (OULAD, UCI Student, Higher Ed, corresponding to Strong, Mostly
Passing, and Partial profiles respectively).

The main audit framework, the four-dimension results, and all numerical
analyses are unchanged. We describe each change below. Specific line
references refer to the fourth-revision manuscript files.

\section*{Response to Reviewer 2}

We are grateful to Reviewer~2 for this final round of careful reading.
The two remaining concerns are addressed below.

\subsection*{Comment~1: Interpretation of the feature-attribution analysis}

\textbf{Reviewer:} The interpretation of the feature-attribution analysis
remains stronger than the evidence supports. Statements describing IID
attribution instability as a ``necessary condition'' for cross-group
fragility, claiming that it ``confirms'' model dependence, or presenting
attribution stability as a ``proxy for dataset readiness'' should be
revised. Since the analysis does not directly assess attribution
stability under group-aware holdout, the findings should be presented as
exploratory associations that are consistent with, but do not validate,
the audit profiles.

\textbf{Response:} We thank the reviewer for this precise and
well-founded observation. The reviewer is correct: because the
feature-attribution analysis is conducted only under IID splits and does
not directly evaluate attribution stability under group-aware holdout,
the previous framing exceeded what the evidence supports. We have
revised the entire feature-attribution section to present the findings
as exploratory associations rather than validation of the audit
profiles.

The specific changes are as follows.

\begin{enumerate}
\item \textbf{``Necessary condition'' removed.}  The previous claim that
  ``Such lower-bound instability is a necessary---though not
  sufficient---condition for cross-group fragility'' has been replaced
  with the explicit scoping statement: ``Because this analysis is
  conducted only under IID splits and does not directly assess
  attribution stability under group-aware holdout, the observed patterns
  provide exploratory associations that are consistent with---but do not
  validate---the audit profiles.''  (behavioraudit.tex, line~325)

\item \textbf{``Confirms'' softened.}  The previous claim that the
  cross-model comparison was ``confirming that attribution rankings from
  fragile datasets are model-dependent as well as split-dependent'' has
  been revised to ``suggesting that attribution rankings from the
  Partial and Mostly Passing datasets are more model-dependent as well
  as split-dependent''.  The verb ``confirm'' has been replaced by
  ``suggest'', and the scope is now restricted to the analysed
  datasets.  (behavioraudit.tex, line~325)

\item \textbf{``Proxy for dataset readiness'' removed.}  The previous
  concluding sentence stated that ``iid-based attribution stability
  provides a proxy for dataset readiness: Strong datasets maintain
  stable attribution patterns across splits, while Fragile datasets
  exhibit split-dependent rankings.''  This claim has been removed.
  The revised conclusion reads: ``OULAD (Strong) maintains stable
  attribution patterns across splits, while UCI Student (Mostly
  Passing) and Higher Ed (Partial) exhibit more split-dependent
  rankings.  These exploratory associations do not, by themselves,
  establish a causal link between IID attribution instability and
  cross-group fragility.''  (behavioraudit.tex, line~325)
\end{enumerate}

We have also re-read the rest of the manuscript (Methods, Results,
Discussion, and Conclusion) and confirmed that other uses of
``confirm'' appear only in the group-holdout and train-test-gap
analyses; we have not modified those instances. The feature-attribution
section now uses only exploratory language throughout.

\subsection*{Comment~2: Scope of the feature-attribution conclusions}

\textbf{Reviewer:} The scope of the feature-attribution conclusions
should be corrected. The analysis includes only OULAD, UCI Student, and
Higher Ed, corresponding respectively to Strong, Mostly Passing, and
Partial profiles; it does not include a dataset formally classified as
Fragile. Therefore, statements referring generally to ``Fragile
datasets'' or claiming a pattern extending from Fragile to Strong
should be replaced with wording restricted to the three datasets
actually analysed.

\textbf{Response:} We agree with the reviewer that the previous
wording generalized the feature-attribution findings beyond the
analysed scope. The feature-attribution analysis includes three
datasets---OULAD (Strong), UCI Student (Mostly Passing), and Higher Ed
(Partial)---and does not include a Fragile-profile dataset. We have
corrected all statements in the feature-attribution section that
referred to ``Fragile datasets'' or implied a Fragile-to-Strong
continuum.

The specific changes are as follows.

\begin{enumerate}
\item \textbf{Opening scope statement.}  The previous sentence stated
  that the three datasets were ``chosen to span Fragile-to-Strong
  profiles''.  The revised text now reads: ``These three datasets were
  chosen to span three of the four readiness profiles (Strong, Mostly
  Passing, and Partial); the Fragile profile (MM-TBA) is not included
  in this analysis because its small sample size and lack of grouping
  metadata make split-level attribution estimates unstable and
  pedagogically uninterpretable.''  (behavioraudit.tex, line~321)

\item \textbf{Profile-specific language.}  In the interpretive
  paragraph, the previous wording ``on Strong data, feature-attribution
  is stable; on Fragile or Partial data, part of the iid signal arises
  from split-specific noise'' has been revised to ``on OULAD (Strong),
  feature-attribution is stable; on UCI Student (Mostly Passing) and
  Higher Ed (Partial), part of the IID signal appears to arise from
  split-specific noise rather than stable underlying structure''.
  (behavioraudit.tex, line~325)

\item \textbf{Cross-model comparison.}  The previous claim that
  ``attribution rankings from fragile datasets are model-dependent as
  well as split-dependent'' has been revised to ``attribution rankings
  from the Partial and Mostly Passing datasets are more model-dependent
  as well as split-dependent''.  (behavioraudit.tex, line~325)

\item \textbf{Concluding summary.}  The previous concluding sentence
  contrasted ``Strong datasets'' with ``Fragile datasets''; the revised
  version explicitly names the three analysed datasets: ``OULAD
  (Strong) maintains stable attribution patterns across splits, while
  UCI Student (Mostly Passing) and Higher Ed (Partial) exhibit more
  split-dependent rankings.''  (behavioraudit.tex, line~325)
\end{enumerate}

In the rest of the manuscript (Methods, Results, Discussion,
Conclusion), the term ``Fragile datasets'' refers to the full
seven-dataset audit, which includes MM-TBA; we have limited the scope
restriction to the feature-attribution section.

\vspace{0.3in}
\noindent We hope these revisions fully address Reviewer~2's remaining
concerns. We sincerely appreciate the reviewer's careful reading over
four rounds of review, which has substantially sharpened the precision
of our claims and the scope of our conclusions. We believe the
manuscript now provides a more rigorous and appropriately scoped
presentation of the audit protocol and its supporting analyses, and we
respectfully submit it for the Editor's final consideration.

\vspace{0.1in}
\noindent Respectfully,\\
Yan Ma \and Lizhuo Zhang

\newpage
\bibliographystyle{unsrtnat}
\bibliography{behavioraudit}

\end{document}
